\section{Plano de Trabalho e Cronograma de atividades}

O plano de trabalho deste projeto contempla a realização de disciplinas obrigatórias exigidas pelo programa de pós-graduação, bem como o desenvolvimento das etapas de pesquisa ao longo dos anos de 2026 e 2027. No primeiro ano, estão previstas atividades acadêmicas relacionadas ao cumprimento de créditos, distribuídas ao longo de todos os trimestres de 2026.

Paralelamente, será conduzida uma revisão bibliográfica inicial nos dois primeiros trimestres de 2026, com o objetivo de fundamentar teoricamente o problema abordado. Essa revisão será continuamente atualizada ao longo do desenvolvimento do projeto, conforme novos avanços relevantes na literatura.

Ainda no primeiro ano, está prevista a realização do Exame de Qualificação de Mestrado (EQM) no terceiro trimestre de 2026, momento no qual serão apresentadas as definições formais do problema, a modelagem proposta e possíveis resultados preliminares.

Adicionalmente, o candidato participará do Programa de Estágio Docente (PED) entre o terceiro e o quarto trimestres de 2026, visando à aquisição de experiência em atividades de ensino e apoio acadêmico.

No que diz respeito ao desenvolvimento científico do projeto, as melhorias nas formulações matemáticas serão realizadas entre o terceiro e o quarto trimestres de 2026, período em que também ocorrerá a consolidação dos modelos propostos.

No segundo ano, está prevista a realização de um estágio de pesquisa no exterior por meio do programa de Bolsa
Estágio de Pesquisa no Exterior (BEPE), a ser desenvolvido ao longo dos dois primeiros trimestres de 2027. Esse estágio tem como objetivo fortalecer a colaboração internacional e aprofundar aspectos teóricos e computacionais do trabalho.

Ainda em 2027, serão desenvolvidas heurísticas baseadas em BRKGA e RKO durante o primeiro semestre, com foco na resolução de instâncias de grande porte. Na sequência, a etapa de experimentação computacional será conduzida no terceiro trimestre, permitindo a avaliação de desempenho dos modelos e métodos propostos.

A análise dos resultados obtidos dará suporte à escrita e submissão de artigos científicos, previstas para o terceiro trimestres de 2027. Por fim, a redação da dissertação será realizada entre o terceiro e o quarto trimestres, culminando na entrega e defesa do trabalho ao final de 2027. A Tabela~\ref{cronograma} demonstra como estas etapas estão distribuídas ao longo do período de desenvolvimento do projeto.

\begin{table}[H]
\caption{Cronograma trimestral do projeto de Mestrado}
\label{cronograma}
\centering
\begin{adjustbox}{width=\textwidth}
\begin{tabular}{|l|cllcllcllcll|cllcllcllcll|}
\hline
                                                   & \multicolumn{12}{c|}{\textbf{2026}}                                                                                                                           & \multicolumn{12}{c|}{\textbf{2027}}                                                                                                                           \\ \hline
                                                   & \multicolumn{3}{c|}{\textbf{Jan-Mar}} & \multicolumn{3}{c|}{\textbf{Abr-Jun}} & \multicolumn{3}{c|}{\textbf{Jul-Set}} & \multicolumn{3}{c|}{\textbf{Out-Dez}} & \multicolumn{3}{c|}{\textbf{Jan-Mar}} & \multicolumn{3}{c|}{\textbf{Abr-Jun}} & \multicolumn{3}{c|}{\textbf{Jul-Set}} & \multicolumn{3}{c|}{\textbf{Out-Dez}} \\ \hline
Disciplinas                                        & \multicolumn{3}{c|}{•}                & \multicolumn{3}{c|}{•}                & \multicolumn{3}{c|}{•}                & \multicolumn{3}{c|}{•}                & \multicolumn{3}{c|}{}                 & \multicolumn{3}{c|}{}                 & \multicolumn{3}{c|}{}                 & \multicolumn{3}{c|}{}                 \\ \hline
EQM                                                & \multicolumn{3}{c|}{}                 & \multicolumn{3}{c|}{}                 & \multicolumn{3}{c|}{•}                & \multicolumn{3}{c|}{}                 & \multicolumn{3}{c|}{}                 & \multicolumn{3}{c|}{}                 & \multicolumn{3}{c|}{}                 & \multicolumn{3}{c|}{}                 \\ \hline
Revisão bibliográfica                              & \multicolumn{3}{c|}{•}                & \multicolumn{3}{c|}{•}                & \multicolumn{3}{c|}{•}                 & \multicolumn{3}{c|}{•}                 & \multicolumn{3}{c|}{•}                 & \multicolumn{3}{c|}{•}                 & \multicolumn{3}{c|}{•}                 & \multicolumn{3}{c|}{•}                 \\ \hline
BEPE                                               & \multicolumn{3}{c|}{}                 & \multicolumn{3}{c|}{}                 & \multicolumn{3}{c|}{}                 & \multicolumn{3}{c|}{}                 & \multicolumn{3}{c|}{•}                & \multicolumn{3}{c|}{•}                & \multicolumn{3}{c|}{}                 & \multicolumn{3}{c|}{}                 \\ \hline
PED                                                & \multicolumn{3}{c|}{}                 & \multicolumn{3}{c|}{}                 & \multicolumn{3}{c|}{•}                & \multicolumn{3}{c|}{•}                & \multicolumn{3}{c|}{}                 & \multicolumn{3}{c|}{}                 & \multicolumn{3}{c|}{}                 & \multicolumn{3}{c|}{}                 \\ \hline
Melhorias nas formulações                          & \multicolumn{3}{c|}{}                 & \multicolumn{3}{c|}{}                 & \multicolumn{3}{c|}{•}                & \multicolumn{3}{c|}{•}                & \multicolumn{3}{c|}{}                 & \multicolumn{3}{c|}{}                 & \multicolumn{3}{c|}{}                 & \multicolumn{3}{c|}{}                 \\ \hline
Desenvolvimento de heurísticas                     & \multicolumn{3}{c|}{}                 & \multicolumn{3}{c|}{}                 & \multicolumn{3}{c|}{}                 & \multicolumn{3}{c|}{}                 & \multicolumn{3}{c|}{•}                & \multicolumn{3}{c|}{•}                & \multicolumn{3}{c|}{}                 & \multicolumn{3}{c|}{}                 \\ \hline
Experimentação computacional & \multicolumn{3}{c|}{}                 & \multicolumn{3}{c|}{}                 & \multicolumn{3}{c|}{}                 & \multicolumn{3}{c|}{}                 & \multicolumn{3}{c|}{}                 & \multicolumn{3}{c|}{}                 & \multicolumn{3}{c|}{•}                & \multicolumn{3}{c|}{}                 \\ \hline
Escrita de artigos cientifícos                     & \multicolumn{3}{c|}{}                 & \multicolumn{3}{c|}{}                 & \multicolumn{3}{c|}{}                 & \multicolumn{3}{c|}{}                 & \multicolumn{3}{c|}{}                 & \multicolumn{3}{c|}{}                & \multicolumn{3}{c|}{•}                & \multicolumn{3}{c|}{}                 \\ \hline
Escrita da dissertação                             & \multicolumn{3}{c|}{}                 & \multicolumn{3}{c|}{}                 & \multicolumn{3}{c|}{}                 & \multicolumn{3}{c|}{}                 & \multicolumn{3}{c|}{}                 & \multicolumn{3}{c|}{}                 & \multicolumn{3}{c|}{•}                & \multicolumn{3}{c|}{•}                \\ \hline
Entrega da dissertação e defesa                    & \multicolumn{3}{c|}{}                 & \multicolumn{3}{c|}{}                 & \multicolumn{3}{c|}{}                 & \multicolumn{3}{c|}{}                 & \multicolumn{3}{c|}{}                 & \multicolumn{3}{c|}{}                 & \multicolumn{3}{c|}{}                 & \multicolumn{3}{c|}{•}                \\ \hline
\end{tabular}
\end{adjustbox}
\end{table}